\documentclass[10pt,a4paper]{amsart}
\usepackage[margin=19mm]{geometry}
\usepackage{amsmath,amssymb,mathtools}
\usepackage[scr=rsfs,cal=euler]{mathalfa}
\usepackage{xfrac}
\usepackage[unicode,hidelinks,bookmarks=false]{hyperref}
\newcommand{\ie}{i.e.\ }
\newcommand{\cf}{cf.\ }
\newcommand{\resp}{respectively\ }
\newcommand{\Subsection}{\subsection}
\providecommand{\nobreakdash}{\nobreak\hbox{-}}
\setlength{\emergencystretch}{2em}
\multlinegap0pt
\usepackage{bm}
\usepackage{bbm}
\usepackage{dsfont}
\usepackage{xspace}
\usepackage{cancel}
\usepackage{mathtools}
\usepackage{amsthm}
\makeatletter
\@ifundefined{proposition}{\newtheorem{proposition}{Proposition}}{}
\makeatother
\title[Copula-Based Modeling of Fractional Inaccuracy: A Unified Fr]{Copula-Based Modeling of Fractional Inaccuracy: A Unified Framework}
\author{Aman Pandey, Chanchal Kundu}
\date{Source: arXiv:2506.19748v1 (CC BY 4.0)}
\begin{document}
\maketitle

\noindent\emph{Source and licence.} Copula-Based Modeling of Fractional Inaccuracy: A Unified Framework. Version arXiv:2506.19748v1 (CC BY 4.0), distributed under the Creative Commons Attribution 4.0 International licence, as recorded in the arXiv licence metadata for this exact version and shown on the abstract page https://arxiv.org/abs/2506.19748v1. Full rights notice: licenses/arxiv-2506.19748-CC-BY-4.0.txt.\par\smallskip

\noindent\textit{Editorial scope.} Counted unit: the Proposition whose source label is \texttt{prop3.6} in the article (source0.tex lines 444--454, source label \texttt{prop3.6}), delivered together with the article's own complete original proof of it (source0.tex lines 456--492, one \texttt{proof} environment of 37 lines). No mathematics is written, repaired or re-derived: statement and proof are the article's own text line for line. Required context retained uncounted: none. Every definition, display and label of those lines is retained unchanged. Omitted from this package and not redistributed: 1--443, 455--455, 493--790. Nothing inside a retained excerpt is omitted or altered: every retained body is delivered exactly as the article's lines are written, so no omission receipt of any kind accompanies this package. Citations: the retained excerpts carry no citation command at all, so no bibliography entry of the article is transcribed and no reference is redistributed. Reference adaptation: none was needed and none was made; every reference inside the delivered unit resolves inside it, so no unresolved reference or citation is produced. Printed slips kept literal: no printed word, symbol, hypothesis or number of the counted statement or of its proof was altered. Layout and class: the article's own class is \texttt{article} with packages including float, multirow, diagbox, amssymb, amsmath, amsfonts, booktabs, bbm, lscape, enumerate, amsthm, fontenc, inputenc, lmodern; the wrapper is the lane's pinned \texttt{amsart} editorial wrapper at 10pt on A4, which restates the theorem-like environments the retained text needs and adds the article's own macro definitions that the retained text needs (none), with extra wrapper packages (bm, bbm, dsfont, xspace, cancel, mathtools). The delivered numbering is the unit's own. Assets: no retained span contains a figure environment, a graphics-inclusion command, a PDF-page inclusion, a table or a TikZ picture; no image file is copied into this package, the package declares no asset dependency and no image file is redistributed with it. Licence: version arXiv:2506.19748v1 of this article is distributed under the Creative Commons Attribution 4.0 International licence, as recorded in the arXiv licence metadata for this exact version and shown on the abstract page https://arxiv.org/abs/2506.19748v1. A full rights notice is delivered at licenses/arxiv-2506.19748-CC-BY-4.0.txt. Characters: every retained excerpt is pure ASCII, so the delivered engine, XeLaTeX with its default Latin Modern text font, reports no missing character.
\begin{proposition}\label{prop3.6}
Suppose that $\mathbb{X},~\mathbb{Y}$ and $\mathbb{Z}$ have copulas $C_\mathbb{X},~C_\mathbb{Y}$ and $C_\mathbb{Z}$, respectively. 
\begin{enumerate}
\item If $\mathbb{Z}\le_{LO}\mathbb{Y},~\mathbb{Z}\le_{LO}\mathbb{X}$ and $Z_i\overset{\mathrm{st}}{=}X_i$,~$i=1,\cdots,n$. 
Then, $CCFI(\mathbb X,\mathbb Z)\leq CCFI(\mathbb X,\mathbb Y)\le CCFI(\mathbb Z,\mathbb Y).$;
\item  If $\mathbb{X}\le_{LO}\mathbb{Z}\le_{LO}\mathbb{Y}$ and $Z_i\overset{\mathrm{st}}{=}X_i$,~$i=1,\cdots,n$. \\Then,
 $CCFI(\mathbb{X},\mathbb{Y})\ge \max\{CCFI(\mathbb{Z},\mathbb{Y})\le CCFI(\mathbb{X},\mathbb{Z})\}$;
\item If $\mathbb{Y}\le_{LO}\mathbb{Z},~\mathbb{X}\le_{LO}\mathbb{Z}$ and$Z_i\overset{\mathrm{st}}{=}X_i$,~$i=1,\cdots,n$. \\Then,
 $CCFI(\mathbb Z,\mathbb Y)\leq CCFI(\mathbb{X},\mathbb{Y})\ge CCFI(\mathbb X,\mathbb Z)$.
\end{enumerate}
\end{proposition}

\begin{proof}
We provide the proof for part $(1)$ only; the remaining parts can be established in a similar manner.
 Since $\mathbb{Z}\le_{LO}\mathbb{X}$ and $Z_i\overset{\mathrm{st}}{=}X_i$,~$i=1,\cdots,n$, we have 
\begin{align}\label{eq3.24}
C_\textbf{Z}(v_1,\cdots,v_n)\ge C_\textbf{X}(v_1,\cdots,v_n),~~u_i\in[0,1].
\end{align}
 From (\ref{eq3.24}), we get
\begin{align*}
& C_\textbf{Z}(v_1,\cdots,v_n)\Big\{-\mathrm{Ln}_\eta C_\textbf{Y}\big(G_1(F^{-1}_1(v_1)),\cdots, G_n(F^{-1}_n(v_n))\big)\Big\}^{\frac{1}{\eta}}\nonumber\\
&\ge C_\textbf{X}(v_1,\cdots,v_n)\Big\{-\mathrm{Ln}_\eta C_\textbf{Y}\big(G_1(F^{-1}_1(v_1)),\cdots, G_n(F^{-1}_n(v_n))\big)\Big\}^{\frac{1}{\eta}}\nonumber\\
\end{align*}
Equivalently,
\begin{align*}
	\int_{\mathbb I}C_\textbf{Z}(v_1,\cdots,v_n)&\Big\{-\mathrm{Ln}_\eta C_\textbf{Y}\big(G_1(F^{-1}_1(v_1)),\cdots, G_n(F^{-1}_n(v_n))\big)\Big\}^{\frac{1}{\eta}}dv_1\cdots dv_n\nonumber\\
	&\ge \ \int_{\mathbb I}C_\textbf{X}(v_1,\cdots,v_n)\Big\{-\mathrm{Ln}_\eta C_\textbf{Y}\big(G_1(F^{-1}_1(v_1)),\cdots, G_n(F^{-1}_n(v_n))\big)\Big\}^{\frac{1}{\eta}}dv_1\cdots dv_n.\nonumber\\
\end{align*}
Thus, we obtain
\begin{align}\label{eq3.25}
	CCFI(\textbf{Z},\textbf{Y})\ge CCFI(\textbf{X},\textbf{Y}).
\end{align}
Further, let $\mathbf Z\leq_{LO}\mathbf Y$. Then
\begin{align*}
	C_{\mathbf Z}(H_1(z_1),...,H_n(z_n))\geq C_{\mathbf Y}(G_1(z_1),...,G_n(z_n))
\end{align*}
Therefore, 
\begin{align*}
	\int_{\mathbb I}C_{\mathbf X}(v_1,...,v_n)&\left(-\mathrm{Ln}_\eta\left(C_{\mathbf Z}(H_1(F_1^{-1}(u_1)),...,H_n(F_n^{-1}(u_n)))\right)\right)^{\frac{1}{\eta}}d\textbf v\leq \\	&\int_{\mathbb I}C_{\mathbf X}(v_1,...,v_n)\left(-\mathrm{Ln}_\eta\left(C_{\mathbf Y}(G_1(F_1^{-1}(u_1)),...,G_n(F_n^{-1}(u_n)))\right)\right)^{\frac{1}{\eta}}d\textbf v,
\end{align*}
which implies 
\begin{align}\label{eq3.26}
	CCFI(\mathbf X,\mathbf Z)\leq CCFI(\mathbf X,\mathbf Y)
\end{align}
After combining (\ref{eq3.25}) and (\ref{eq3.26}), we obtain
\begin{align*}
CCFI(\mathbf X,\mathbf Z)\leq CCFI(\textbf X,\textbf Y)\le CCFI(\mathbf Z,\mathbf Y).
\end{align*}
\end{proof}

\end{document}
